\section{Representações de Grafos}

\begin{frame}{Lista de Adjacência}
    Um grafo $G=(V,E)$ é representado por um arranjo de listas. Para cada vértice $u \in V$, há uma lista que contém todos os seus vizinhos.
    
    $$ \text{Adj}(u) = \{ v \in V \mid (u, v) \in E \} $$
    
    \pause % <-- O slide para aqui até você avançar
    
    \textbf{Pontos Fortes}
    \begin{itemize}
        \item<+-> Eficiência de Espaço: Ideal para grafos com poucas arestas (esparsos). O espaço usado é proporcional ao número de vértices e arestas: $O(|V|+|E|)$.
        \item<+-> Busca de Vizinhos Rápida: Listar todos os vizinhos de um vértice é muito eficiente.
    \end{itemize}

    \pause % <-- O slide para aqui até você avançar
    
    \textbf{Pontos Fracos}
    \begin{itemize}
        \item<+-> Verificação de Aresta Lenta: Checar se a aresta (u,v) existe pode exigir a varredura de toda a lista de vizinhos de u.
    \end{itemize}
\end{frame}

\begin{frame}{Matriz de Adjacência}
    Um grafo com $n$ vértices é representado por uma matriz quadrada $M$ de dimensão $n \times n$.
    
    $$
    M_{uv} =
    \begin{cases}
        1 & \text{se existe a aresta } (u, v), \\
        0 & \text{caso contrário.}
    \end{cases}
    $$
    
    \pause % Primeira pausa
    
    \textbf{Pontos Fortes}
    \begin{itemize}
        \item<+-> Verificação de Aresta Instantânea: Saber se a aresta (u,v) existe é uma operação de tempo constante: $O(1)$.
        \item<+-> Adição/Remoção Rápida de Arestas: Alterar o valor na matriz de 0 para 1 (ou vice-versa) é muito rápido.
    \end{itemize}    

    \pause % <-- O slide para aqui até você avançar
    
    \textbf{Pontos Fracos}
    \begin{itemize}
        \item<+-> Alto Custo de Memória: O espaço é sempre $O(|V|^2)$, mesmo para grafos com poucas arestas, o que gera desperdício.
    \end{itemize}
\end{frame}

\begin{frame}{Exemplo: Sugestão de Amigos em Rede Social}
    \textbf{O Problema:} Para um usuário, encontrar seus "amigos de amigos" para sugerir novas conexões.
    
    \pause % Pausa 1
    
    \begin{itemize}
        \item \textbf{Modelo:} Um grafo com 1 milhão de vértices (usuários).
        \item \textbf{Característica:} Grafo \textit{esparso}, onde cada usuário tem em média 200 amigos (arestas).
    \end{itemize}
    
    \pause % Pausa 2
    
    \textbf{Abordagem 1: Usando Lista de Adjacências}
    \begin{itemize}
        \item<+-> \textbf{Passo 1:} Acessar a lista do usuário inicial. Custo: ler 200 nomes.
        \item<+-> \textbf{Passo 2:} Para cada um desses 200 amigos, acessar suas respectivas listas (com ~200 nomes cada).
    \end{itemize}
    
    \pause % Pausa 3
    
    \begin{block}{Custo Computacional}
        O número de operações é proporcional às conexões reais:
        $$ \approx 200 \times 200 = \textbf{40.000 operações (rápido e eficiente)} $$
    \end{block}

\end{frame}

\begin{frame}{Análise Comparativa e Conclusão}
    
    \textbf{Abordagem 2: Usando Matriz de Adjacências}
    
    \begin{itemize}
        \item<+-> \textbf{Custo de Memória:} Exigiria uma matriz de $10^6 \times 10^6$, ou seja, $10^{12}$ células. \textbf{Totalmente inviável!}
        \item<+-> \textbf{Custo de Computação:} Para achar os amigos de alguém, é preciso varrer 1 milhão de colunas.
            $$ \approx 200 \times 1.000.000 = \textbf{200 milhões de operações} $$
    \end{itemize}

    \pause % Pausa 1
    
    \begin{block}{Tabela Comparativa}
    \centering
    \begin{tabular}{r c c}
        \hline
        \textbf{Critério} & \textbf{Lista de Adjacências} & \textbf{Matriz de Adjacências} \\
        \hline
        Memória & Mínima ($O(|V|+|E|)$) & Gigantesca ($O(|V|^2)$) \\
        Velocidade & Quase Instantâneo & Extremamente Lento \\
        \hline
    \end{tabular}
    \end{block}
    
    \pause % Pausa 2

    \textbf{Conclusão:} Para grafos esparsos do mundo real, a \textbf{Lista de Adjacências} é a única escolha computacionalmente viável e eficiente.

\end{frame}

\begin{frame}{Exemplo de IA: Análise de Relações em Textos}
    \textbf{O Problema:} Um modelo de IA (como os que equipam a busca do Google ou o ChatGPT) precisa aprender a relação semântica entre as palavras mais importantes de um idioma.

    \pause % Pausa 1
    
    \begin{itemize}
        \item \textbf{Modelo:} Um grafo com os \textbf{5.000 termos} mais comuns de um vocabulário.
        \item \textbf{Característica:} Grafo \textit{denso}. Palavras como "governo" ou "economia" se conectam com milhares de outras.
        \item \textbf{Operação Crítica:} Durante o treinamento, o modelo precisa verificar milhões de vezes e de forma \textbf{instantânea} se dois termos, como ('ciência', 'dados'), possuem uma conexão direta.
    \end{itemize}
    
    \pause % Pausa 2
    
    \textbf{Abordagem 1: Usando Matriz de Adjacências}
    \begin{itemize}
        \item<+-> \textbf{Estrutura:} Uma matriz $5000 \times 5000$. O espaço de memória é alto, mas fixo e perfeitamente gerenciável para este porte.
        \item<+-> \textbf{Verificação de Aresta:} Para checar a conexão ('ciência', 'dados'), o acesso é direto à célula \texttt{M[id\_ciencia][id\_dados]}.
    \end{itemize}
\end{frame}        

\begin{frame}{Exemplo de IA: Análise de Relações em Textos}
    \begin{block}{Custo Computacional}
        A verificação da existência de uma aresta é em tempo constante:
        $$ O(1) $$        
    \end{block}
\end{frame}

\begin{frame}{Análise Comparativa e Conclusão (Grafos Densos)}
    \textbf{Abordagem 2: Usando Lista de Adjacências (Neste Cenário)}
    \begin{itemize}
        \item \textbf{Verificação de Aresta:} Para checar a conexão ('ciência', 'dados'), o algoritmo precisaria:
        \begin{enumerate}
            \item Ir para a lista de vizinhos da palavra 'ciência'.
            \item Percorrer \textbf{toda a lista} (que pode ter milhares de palavras) para procurar por 'dados'.
        \end{enumerate}
        \item<+-> \textbf{Custo:} $O(\text{grau}(\text{ciência}))$. Em um grafo denso, isso significa \textbf{milhares de operações para uma única verificação}, o que tornaria o treinamento da IA extremamente lento.
    \end{itemize}
    
    \pause % Pausa 1
    
    \begin{block}{Tabela Comparativa (Para este Cenário de IA)}
        \centering
        \begin{tabular}{r c c}
            \hline
            \textbf{Critério} & \textbf{Matriz de Adjacências} & \textbf{Lista de Adjacências} \\
            \hline
            Verif. de Aresta & \textbf{Instantâneo ($O(1)$)} & Lento ($O(grau)$) \\
            Uso em IA & \textbf{Ideal para treino} & Ineficiente para treino \\
            \hline
        \end{tabular}
    \end{block}
    
    \pause % Pausa 2
    
    \textbf{Conclusão:} Para grafos \textbf{densos} e checagens de arestas em \textbf{$O(1)$}, a \textbf{Matriz de Adjacências} é a mais eficiente.
\end{frame}